\documentclass[11pt,a4paper]{article}
\usepackage{graphicx}
\usepackage[utf8]{inputenc}
\usepackage[T1]{fontenc}
\usepackage[english]{babel}
\usepackage{lmodern,microtype}
\usepackage[a4paper,margin=2.4cm,headheight=15pt]{geometry}
\usepackage{setspace,parskip,enumitem,booktabs,longtable,xcolor,hyperref,fancyhdr,titlesec,listings}
\definecolor{bleu}{RGB}{26,65,110}
\definecolor{lien}{RGB}{20,75,140}
\hypersetup{colorlinks=true,linkcolor=bleu,urlcolor=lien,pdftitle={Second Progress Report -- Agent Capability Graph}}
\setstretch{1.1}
\setlist[itemize]{leftmargin=1.5em,itemsep=0.25em}
\setlist[enumerate]{leftmargin=1.7em,itemsep=0.35em}
\titleformat{\section}{\Large\bfseries\color{bleu}}{\thesection}{0.7em}{}
\titleformat{\subsection}{\large\bfseries\color{bleu}}{\thesubsection}{0.7em}{}
\pagestyle{fancy}\fancyhf{}\lhead{\small Second Progress Report}\rhead{\small Agent Capability Graph}\cfoot{\thepage}
\lstset{basicstyle=\ttfamily\small,breaklines=true,frame=single}

\begin{document}
\hypersetup{pageanchor=false}
\begin{titlepage}
\centering\vspace*{1.7cm}
{\Large\bfseries MY 2nd PROGRESS REPORT \par}
\vspace{1.2cm}
{\Huge\bfseries\color{bleu} Agent Capability Graph\par}
\vspace{0.5cm}
{\LARGE Building an RDF Knowledge Graph to Discover MCP Servers and Agent Skills\par}
\vspace{1.4cm}
\begin{tabular}{@{}ll@{}}
\textbf{Main Supervisor:} & Raphael Troncy\\
\textbf{Secondary Supervisor:} & Thibault Ehrhart\\
\textbf{Document Type:} & Progress Report\\
\textbf{Phase:} & Dataset and RDF Pipeline Prototype\\
\end{tabular}

\vfill

\vspace{1cm}{\large \today\par}
\end{titlepage}
\hypersetup{pageanchor=true}

\section{Project goal}

The project explores how an RDF knowledge graph can help users discover and
compare MCP servers, their tools, and Agent Skills. The graph is intended to
connect source-backed descriptions with normalized capability information,
while keeping track of where each metadata value came from. The current work is
a prototype and a curated sample, rather than a general-purpose web crawler.

\section{Dataset collection}

The expanded sample contains ten MCP servers, forty-three selected documented
tools, and ten Agent Skills. The records come from nine GitHub repositories.
The MCP sample includes the official MCP server repository and projects from
GitHub, Brave, Cloudflare, and AWS. The Skills sample includes repositories
from Anthropic, Vercel, Microsoft, and Google.

\begin{center}
\begin{tabular}{@{}lr@{}}
\toprule
\textbf{Collected item} & \textbf{Count}\\
\midrule
MCP servers & 10\\
Documented tool samples linked to servers & 43\\
Agent Skills & 10\\
Distinct source repositories & 9\\
Downloaded Skill package archives & 10\\
\bottomrule
\end{tabular}
\end{center}

The collector records repository revisions, source URLs, field values, and
evidence statuses in SQLite and exports a CSV for review. It saves the reviewed
MCP README files and downloads each selected Skill directory as a ZIP package.
The downloaded files are not executed. The tool list is a documented sample
transcribed from the selected source material; it is not guaranteed to include
every tool available from a server at runtime.

\section{Ontology and graph model}

The first ontology version is written in Turtle in
\texttt{docs/acg\_ontology.ttl}. It defines classes for MCP servers, MCP tools,
Agent Skills, capabilities, source documents, metadata observations, and
capability assertions. It also defines relationships such as
\texttt{acg:exposesTool}, which connects a server to its tools, and provenance
properties that record reviewed sources and observation status.

The ontology reuses terms from Schema.org, SKOS, PROV-O, RDF, RDFS, and OWL where
appropriate. Its namespace, \texttt{urn:acg:ontology:}, is still provisional.
The model distinguishes source descriptions from curator-assigned capability
labels and preserves source evidence for later review.

\section{SQLite-to-RDF pipeline}

The current workflow has two main steps:

\begin{enumerate}
  \item The expanded collector harvests the curated sample and writes the
  SQLite database and CSV under the interim-data directory.
  \item The RDF generator maps the records to the ontology and writes the
  harvested Turtle graph under the processed-data directory.
\end{enumerate}

The generated graph currently contains 2,210 RDF triples representing 63
harvested entities: 10 servers, 43 tools, and 10 Skills. Each collected field is
retained as a source-backed metadata observation. The conversion also links each
tool to its parent server and emits direct RDF properties for selected fields.

The generator parses both the ontology and the output graph, and checks that
used ACG classes and properties are declared, that tools have server links, and
that metadata observations are complete. These application-level checks are
useful for the current pipeline, but they are not a full OWL reasoning step or
SHACL validation. SHACL remains a possible later validation step, not a current
project requirement.

\section{Project organization}

The code is organized in \texttt{src/acg/}. Small root-level launchers preserve
the commands used during earlier development. Data is separated into
\texttt{data/sources/} for maintained inputs,
\texttt{data/raw/downloads/} for retrieved source packages,
\texttt{data/interim/} for the SQLite and CSV collection outputs, and
\texttt{data/processed/} for the generated RDF graph. Project reports and
ontology documents are kept under \texttt{docs/}. A \texttt{.gitignore} file
excludes local Python environments, cache files, and downloaded packages.

\section{Current limitations and next steps}

The current sample and pipeline provide a working base for evaluating the
ontology, but several steps remain:

\begin{itemize}
  \item Refine competency questions from the perspective of graph users and
  state the expected answer for each question.
  \item Test whether the generated graph can answer those questions using
  SPARQL queries.
  \item Improve direct RDF mappings for useful metadata such as repository
  identity, Skill compatibility, and package archive location. These values are
  currently preserved as observations when no dedicated mapping is applied.
  \item Review the sample and its source evidence, and distinguish clearly
  between source-declared facts and curator interpretations.
  \item Consider SHACL shapes after the ontology and data fields stabilize.
\end{itemize}

\section{Conclusion}

The project has progressed from a small exploratory JSON-LD example to a
multi-source dataset and a reproducible SQLite-to-RDF prototype. The next
milestone is to connect the ontology and graph to concrete user questions, then
refine the model and mappings based on the answers those questions require.

\end{document}
